% TOP_PROBLEM: {"id": 3142, "title": "Chromatic number of $\\frac{3}{3}$-power of graph", "category_id": 3, "category": {"id": 3, "name": "graph_theory", "display_name": "Graph Theory", "description": "Problems involving graphs, networks, and their properties.", "slug": "graph-theory", "order_index": 3, "created_at": "2026-07-31T15:26:25.671Z"}, "status": "open", "problem_number": "OPG-60055", "set_id": 12}
% FIRST_POSTED: 2026-10-01
% ATTEMPT_STATUS: unresolved
% UnsolvedMath ID 3142; source catalogue code OPG-60055
% !TeX program = lualatex
% Research attempt by GPT-6 Astra Ultra; no claim of a new complete solution.
\documentclass[11pt,a4paper]{article}
\usepackage[margin=25mm]{geometry}
\usepackage{fontspec}
\setmainfont{lmroman10-regular.otf}[BoldFont=lmroman10-bold.otf,
 ItalicFont=lmroman10-italic.otf,BoldItalicFont=lmroman10-bolditalic.otf]
\usepackage{amsmath,amssymb,mathtools}
\usepackage{unicode-math}
\setmathfont{latinmodern-math.otf}
\AtBeginDocument{\let\setminus\smallsetminus}
\usepackage{xurl,hyperref}
\hypersetup{colorlinks=true,urlcolor=blue}
\setcounter{secnumdepth}{0}
\setlength{\parindent}{0pt}
\setlength{\parskip}{5pt}
\setlength{\emergencystretch}{3em}
\begin{document}
\section*{Chromatic number of $\frac{3}{3}$-power of graph}\label{3142}
{\small UnsolvedMath ID 3142 · OPG-60055 · Graph Theory · OpenGarden}
\subsection{Definitions and mathematical statement}
All graphs here are finite, simple, and undirected. For an integer
$n\geq1$, the $n$-subdivision $G^{1/n}$ is obtained by replacing
each edge $uv$ of $G$ by a path of exactly $n$ edges from $u$ to
$v$. The $n-1$ internal vertices on each such path are new, and
internal vertices belonging to different original edges are
distinct. For an integer $m\geq1$, the $m$th power of a graph $H$
joins two distinct vertices when their distance in $H$ is at most
$m$. Vertices in different components have infinite distance.

Define the fractional-power notation by
\[
 G^{m/n}=(G^{1/n})^m.
\]
In particular, $G^{3/3}$ means the third power of the graph obtained
by subdividing every edge twice. The fraction in this notation is
not reduced: $G^{3/3}$ need not be $G$.

Write $\Delta(G)$ for the maximum vertex degree of the original
graph. A proper $q$-colouring assigns one of $q$ colours to each
vertex, assigning different colours at adjacent vertices; $\chi$
is the least possible $q$. The conjecture is
\[
 \chi(G^{3/3})\leq2\Delta(G)+1
 \qquad\text{for every }G\text{ with }\Delta(G)\geq2.
\]
Consequently the desired colours must separate every pair within
three steps in the subdivided graph, including both original and
new vertices. The degree in the bound is measured before either
operation.
\subsection{Short English statement}
After replacing each edge by a three-edge path, can all vertices within distance three be distinguished using at most twice the original maximum degree plus one colours?
\subsection{Research attempt}
\paragraph{Scope and process.}
This is a GPT-6 Astra Ultra research attempt dated 1 October 2026, using
primary-source browsing and elementary mathematical checks. The canonical
definition above is retained. Results below are restricted results or
reductions, not a claim of a new solution to the entire catalogue question.
No formal proof assistant or external mathematical referee checked this text.


\paragraph{Literature and scope.}
Mozafari-Nia and Iradmusa [R1] prove the conjectured bound in the
subquartic case and state the general $2\Delta+1$ question.
We independently construct colourings for all graphs of maximum
degree two and for all trees. The tree argument uses the actual
three-edge subdivision; reducing the notation $3/3$ to $1$ would
change the graph and invalidate these calculations.

\paragraph{Cycles: a five-colour construction.}
Subdividing $C_m$ produces the cycle $C_{3m}$, with $m\ge3$.
Let $N=3m$. Every such $N\ge9$ can be written $4a+5b$ with
nonnegative integers $a,b$: $9=4+5$, and all integers at least
$12$ have this form, since $12,13,14,15$ do and adding four
preserves it. Arrange around the cycle a concatenation of $a$
blocks $(1,2,3,4)$ and $b$ blocks $(1,2,3,4,5)$, in any order.

No two equal colours have cyclic distance at most three. For
colours $1,2,3,4$, consecutive occurrences are in consecutive
blocks and separated by the preceding block length, four or
five. For colour five, consecutive occurrences are separated by
at least five; if it occurs just once there is no pair to check.
The same reasoning applies across the closing boundary. Hence
the colouring is proper for $C_N^3$ and uses at most five colours.
For $m=3$, any independent set of $C_9^3$ has at most two vertices:
three selected vertices would require three cyclic gaps each at
least four, totaling at least twelve. Thus $C_9^3$ needs at
least $\lceil9/2\rceil=5$ colours, so the bound is attained.

\paragraph{Paths and disconnected graphs of degree two.}
A path with $m$ edges subdivides to a path with $3m$ edges.
Colour its successive vertices periodically $1,2,3,4$. Vertices
at distance one, two or three have different colours, giving
a four-colouring of its third power. Isolated vertices need
one colour. Every finite graph of maximum degree at most two
is a disjoint union of these paths and cycles; separate components
have no finite-distance adjacency, so palettes may be reused.
This proves the target bound $2\Delta+1=5$ for $\Delta=2$.

\paragraph{A stronger bound for every tree.}
Let $H$ be any tree, rooted at a vertex, and order its vertices
by nonincreasing depth. The neighbours of a vertex $v$ in $H^3$
which occur later in this order form a clique. To check this,
write $p$ for its parent. Such neighbours consist only of ancestors
within three steps, siblings of $v$, and children of its
grandparent of depth one less than $v$; some same-depth siblings
may have been removed already. No later descendant of $v$ is
possible. Every pair in this list is at distance at most three,
as is seen by their paths through the parent and grandparent.
The same description truncates harmlessly near the root.
Thus colouring in reverse order needs at most the largest clique
size: each vertex sees a clique of previously coloured neighbours.

A set of vertices pairwise within distance three in a tree lies
within the closed neighbourhood of some edge, or within one
closed vertex-neighbourhood. Indeed its minimal connecting subtree
has diameter at most three, with its center an edge or vertex.
For a central edge $uv$ the union $N_H[u]\cup N_H[v]$ has
$d_H(u)+d_H(v)$ vertices. If now $H=G^{1/3}$ for a tree $G$
of maximum degree $\Delta\ge2$, every edge of $H$ either joins
an original vertex to a degree-two subdivision vertex, or joins
two degree-two vertices. The largest possible clique size is
therefore at most $\Delta+2$. We obtain the stronger restricted
bound
\[
 \chi(G^{3/3})\le\Delta+2\le2\Delta+1
 \qquad(G\text{ a tree},\ \Delta\ge2).
\]

\paragraph{Verification and remaining obstruction.}
The script
\texttt{scripts/check\_attempt\_batch02\_algebra\_2026\_10\_01.py}
checks the cyclic block colourings for $9\le N\le300$ divisible
by three. The proofs, rather than that cutoff, cover all sizes.
The tree elimination argument uses unique paths. In a subdivided
graph with cycles, vertices later in the order can acquire short
routes around a cycle and need not form a clique. Consequently
the tree bound does not extend by choosing a spanning tree:
additional original edges introduce additional subdivision vertices
and new distance-three conflicts.

\paragraph{Final status.}
\textbf{UNRESOLVED} for arbitrary maximum degree and cyclic graphs.
All degree-two instances and all trees are settled by explicit
colourings or a proved elimination rule, with a sharp degree-two
example and a precise failure of the spanning-tree shortcut.

\subsection{Sources}
\begin{itemize}
\item[S1] ULAM.AI and the UnsolvedMath Contributors, supplied problems.json, record 3142 (OPG-60055), OpenGarden collection. Catalogue identity and source transcription; CC BY 4.0.
\url{https://huggingface.co/datasets/ulamai/UnsolvedMath}.
\item[S2] Open Problem Garden, Chromatic number of $\frac{3}{3}$-power of graph. Original problem and discussion; the mathematical formulation below is expanded from the supplied source record.
\url{http://www.openproblemgarden.org/op/chromatic_number_of_frac_3_3_power_of_graph}.

\item[R1] Mahsa Mozafari-Nia and Moharram N. Iradmusa,
\emph{A note on coloring of $3/3$-power of subquartic graphs},
Australasian Journal of Combinatorics 79(3) (2021), 454--460.
Primary paper, definition and main theorem checked 1 October 2026.
\url{https://ajc.maths.uq.edu.au/pdf/79/ajc_v79_p454.pdf}.

\end{itemize}
\end{document}
